
\documentclass[a4paper,12pt]{article}

\input{glyphtounicode}
\pdfgentounicode=1

\usepackage{mathptmx}
\usepackage{fancyhdr}
\usepackage{hyperref}
\usepackage{footmisc}
\usepackage[english]{babel}
\usepackage[utf8]{inputenc}
\usepackage[version=4]{mhchem}
\usepackage[usenames,dvipsnames]{color}
\usepackage[
  a4paper,
  top    = 2.00cm,
  bottom = 2.00cm,
  left   = 1.70cm,
  right  = 1.70cm
]{geometry}

\renewcommand{\baselinestretch}{1.25}
\renewcommand{\headrulewidth}{0pt}
\renewcommand{\footrulewidth}{0pt}
\renewcommand\footnotelayout{\fontsize{10}{10}\selectfont} 

\hypersetup{
    colorlinks=true,
    linkcolor=RoyalBlue,
    urlcolor=RoyalBlue,
    citecolor=RoyalBlue,
    anchorcolor=RoyalBlue
}

\urlstyle{same}
\pagestyle{fancy}
\fancyhf{}
% \fancyfoot[R]{\thepage}

\begin{document}

\begin{center}
  {\Huge \scshape {\fontsize{25}{30}\selectfont{Reza}} {\fontsize{25}{30}\selectfont{\textbf{Aliasgari Renani}}}} \\ 
  {\fontsize{12}{12}\selectfont 25/12/2025} $|$
  {\fontsize{12}{12}\selectfont Motivation Letter} $|$
  {\fontsize{12}{12}\selectfont TUDelft} $|$
  {\fontsize{12}{12}\selectfont PhD Position (2955)} $|$
  {\fontsize{12}{12}\selectfont Integrated Circuit Design}
\end{center}

\vspace{-20pt}
\noindent\rule{\textwidth}{0.5pt}

My research objective is to advance low-power mixed-signal integrated circuits by creating flexible ultrasound front-ends and signal chains that interface efficiently with distributed MEMS transducer arrays. The PhD position in Integrated Circuit Design for Wearable Ultrasound at TU Delft is an excellent opportunity to pursue this aim by developing ASICs that drive and process signals from these arrays, with emphasis on circuit architecture, integration, and measurement. I am now completing the final year of my Master's degree in Plasma Physics at Moscow Institute of Physics and Technology, where my research on semiconductor devices exposed to radiation and FPGAs have strengthened my interest in reliable circuit behavior and digital signal processing systems. 

\vspace{10pt}
I joined Dr.\ Koveshnikov's laboratory in the Russian Academy of Sciences over two years ago to study radiation effects on microelectronic structures, focusing on electrical characterization of semiconductors and \ce{SiO2}-based MOS devices. We studied the impact of electron and gamma irradiation and hydrogen plasma treatment on MOS devices, identified traps by location (oxide, semiconductor, interface) and carrier type (electrons, holes). This work has resulted in a first-author publication\footnote{R. Aliasgari Renani, O.A. Soltanovich, M.A. Knyazev, S.V. Koveshnikov, Investigation of low energy electron irradiated \ce{SiO2} based MOS devices by C-V and TSC techniques, Russian Microelectronics, 2023} and provides insight into device stability under bias stress, radiation immunity, and oxide trap density at the dielectric-semiconductor interface. I also developed optimized MATLAB applications to automate the experimental characterization techniques, which reduced run-times and created measurement pipelines. 

\vspace{10pt}
After my undergraduate degree, I joined the System-on-Chip Development Laboratory and the Laboratory of Plasma Systems under the supervision of Dr.\ Vasiliev to develop FPGA-based systems and to study radiation and plasma effects on these structures. I implemented video-processing algorithms by manual Verilog coding and by generating HDL from Simulink models. I verified functionality of the algorithms using testbenches and Python automation scripts. In the Plasma Lab, we ran experiments where an FPGA with an attached camera and a synthesized image-processing design was placed in an electron-beam plasma system with low-pressure oxygen and exposed to electron irradiation and X-rays generated from a tungsten target. The FPGA output signal was monitored while thermal cycling was applied and preliminary functionality tests were performed. 

\vspace{10pt}
My experience with radiation effects in microelectronic structures, experimental automation, and FPGA based system development motivates me to pursue a PhD where low-power analog and mixed-signal ASICs are designed, prototyped, and validated for wearable ultrasound systems. In this project, I aim to contribute to driver and receive chains for flexible MEMS ultrasound arrays, including circuit architectures, system integration on flexible substrates, and experimental characterization of fabricated silicon and PCB prototypes. I am particularly interested in combining circuit design with FPGA-based control and automated lab measurements to enable robust wearable ultrasound patches for continuous monitoring. It would be an honor to study and work in the Smart Ultrasound Group at the Electronic Instrumentation Laboratory under the supervision of Dr.\ Pertijs at TUDelft.


\end{document}
